\documentclass[compress]{beamer}
% Nord theme: Copyright 2020 by Junwei Wang <i.junwei.wang@gmail.com>

\usepackage{preamb-japanese}
\usepackage{preamb}

\title{Beamer で Nord というテーマを使ってみる}
\subtitle{わくわくしちゃうよね}
\author{俺}
\institute{京都大学}
\date{\today}

\begin{document}

\begin{frame}[plain,noframenumbering]
	\maketitle
\end{frame}


\section{Appearance 外観みたいなものかな？知らんけど。ビーマーノード}

\begin{frame}[fragile]
	\frametitle{Test lstlistings}
	\begin{lstlisting}[language=C++]
#include <bits/stdc++.h>
using namespace std;

int main() {
    int a;
    cout << "Hello world!" << endl;
    // 色は匂へど 散りぬるを
    cout << "我が世誰ぞ 常ならむ" << endl;
    // 有為の奥山 今日越えて
    // 浅き夢見し 酔ひもせず
    // The quick brown fox jumps over the lazy dog.
    cin << a;
    a *= 2;
    cout << a << endl;
    return 0;
}
\end{lstlisting}
\end{frame}


\begin{frame}[fragile]
	\frametitle{Usage 使う}
	Simply include the following code in your preamble:

	これを入れるだけでOK:

	\begin{lstlisting}[language=tex]
\usetheme{Nord}
\end{lstlisting}

	\bigskip

	By default, the appearance is in dark theme, however you can actively choose a either a light or a
	dark theme.

	デフォルトでは黒い背景になりますが、白い背景も選択できます。文字もじモジABCdefご機嫌よう。
    行距line spacing看起來還可以嗎？應該可吧。

	\begin{lstlisting}[language=tex]
\usetheme[style=light]{Nord}
\usetheme[style=dark]{Nord}
\end{lstlisting}

\end{frame}

\subsection{Colors}

\begin{frame}{Defined Colors}{This is a subtitle}
	\begin{description}[Snow Storm]
		\item[Polar Night]
		      \textcolor{NordDarkBlack}{NordDarkBlack} \quad \textcolor{NordBlack}{NordBlack}\\
		      \textcolor{NordMediumBlack}{NordMediumBlack} \quad \textcolor{NordBrightBlack}{NordBrightBlack}
		\item[Snow Storm]
		      \textcolor{NordWhite}{NordWhite} \quad \textcolor{NordBrighterWhite}{NordBrightestWhite}\\
		      \textcolor{NordBrightestWhite}{NordBrightestWhite}
		\item[Forest]
		      \textcolor{NordCyan}{NordCyan} \quad \textcolor{NordBrightCyan}{NordBrightCyan}\\
		      \textcolor{NordBlue}{NordBlue} \quad \textcolor{NordBrightBlue}{NordBrightBlue}
		\item[Aurora]
		      \textcolor{NordRed}{NordRed} \quad \textcolor{NordOrange}{NordOrange} \\
		      \textcolor{NordYellow}{NordYellow} \quad \textcolor{NordGreen}{NordGreen} \\
		      \textcolor{NordMagenta}{NordMagenta}
	\end{description}
\end{frame}

\subsection{Fonts}

\begin{frame}[fragile]{Recommended Free Fonts}
	\begin{description}[Selected Fonts]
		\item[Font選択] recommended for this theme\\
            \begin{lstlisting}[language=tex]
% これは嘘です
\setmainfont{Yanone Kaffeesatz}
\setsansfont{Andika New Basic}
\setmonofont{DejaVu Sans Mono}
\end{lstlisting}
		\item[ダウンロード] {\small \url{https://www.fontsquirrel.com/}}
		\item[インストール] {\small \url{https://www.google.com/get/noto/help/install/}}
		\item[コンパイル] compile with \XeLaTeX~to use system-wide fonts
	\end{description}

\end{frame}


\subsection{Blocks}

\begin{frame}
	\frametitle{Blocks}
	\begin{block}{This is a Block}
		\[
			a^2 + b^2 = c^2
		\]
	\end{block}
	\begin{exampleblock}{This is an Example Block}
		\[
			E = m \cdot c^{2}
		\]
	\end{exampleblock}
	\begin{alertblock}{注意ブロックです}
		\[
			e^{i\pi} + 1 = 0
		\]
	\end{alertblock}

	\centering
	\begin{minipage}{1.0\linewidth}
		\begin{block}{Horizontally-Aligned Block}
			\[
				\log xy = \log x + \log y
			\]
		\end{block}
	\end{minipage}
\end{frame}

\subsection{Items}

\begin{frame}{Items}
	Itemize
	\begin{itemize}
		\item item 1
		\item item 2
	\end{itemize}

	\bigskip

	Enumerate
	\begin{enumerate}
		\item item 1
		\item item 2
		\item 項目３
	\end{enumerate}
\end{frame}

\subsection{Figures}

\begin{frame}{Figures}
	\begin{figure}
		\centering
		\begin{tikzpicture}
			\draw [help lines,NordMagenta,very thick] (0,0) grid (5,4);
		\end{tikzpicture}
		\caption{Credits to Ti\textit{k}Z}
	\end{figure}
\end{frame}

\end{document}

%%% Local Variables:
%%% mode: latex
%%% TeX-master: t
%%% TeX-engine: xetex
%%% End:
